\documentclass[11pt,a4paper]{article}
\usepackage[T1]{fontenc}
\usepackage{lmodern,amsmath,amssymb,amsthm,mathtools,microtype}
\usepackage[margin=27mm]{geometry}
\usepackage[colorlinks=true,linkcolor=blue,citecolor=blue,urlcolor=blue]{hyperref}
\setlength{\emergencystretch}{2em}
\newtheorem{theorem}{Theorem}[section]
\newtheorem{lemma}[theorem]{Lemma}
\newcommand{\Q}{\mathbb Q}\newcommand{\Z}{\mathbb Z}
\DeclareMathOperator{\Tr}{Tr}\DeclareMathOperator{\Nm}{N}
\title{Thirteen-variable conditional sources with intrinsic congruent-number poles}
\author{}\date{17 November 2026}
\begin{document}\maketitle
\begin{abstract}
We replace both pole blocks in the conditional elliptic sources by a
degree-nine radicand of coefficient height four. Multiplying the
pole inputs by seven preserves every consumed odd contact valuation. Ordered
completeness, the required prepared input conditions and every
finite-algebra character guard are verified using Burungale's
simultaneous-twist theorem. The total count remains thirteen variables.
No numerical determinant-degree bound or nine-variable definition is asserted.
\end{abstract}

\section{Specification and premises}
Retain the curve, three point roles $A,H,U$, square-abscissa charts,
subgroup assertion $2E(\Q(\sqrt2))=\Z P$, formal comparison and
trace-height theorem of \cite{Conditional,Reweighted,Unit}. Retain the
simultaneous prime-pattern input in its stated rational-contact form,
including conjugate reduction orders, finite avoidance, denominator
divisibility and real approximation. The former pole-realization
premise is replaced by Burungale's Theorem B in the precise form
established sufficient in \cite{Intrinsic}.

Choose $(\kappa,e)=(121/100,6)$ or $(16/15,4)$ and assume the same
weak rational $abc$ inequality as in \cite{Conditional}. For a fixed
odd $B\ge3$, retain
\[
 \beta=X^2+B,\quad T=\Tr x(A),\quad a_0=1,\quad
 a_1=\beta^e+T^e,\quad t_j=a_j\lambda\quad(j=0,1),\qquad\lambda\ne0.
\]
The six tests are unchanged from the fast conditional detector:
\[
 \begin{gathered}
 f=1/x,\quad z=-x/y,\quad r=z(U)/z(H),\quad
 s=z(U-A)/z(H-P),\\
 C_1=f(H),\quad C_2=f(U),\quad C_3=\overline{f(H-P)},\quad
 C_4=\overline{f(U-A)},\\
 C_5=r-\bar s,\qquad C_6=r+\bar s-2X.
 \end{gathered}
\]
At each argument $t=t_j$ put
\begin{equation}\label{eq:detector}
 \begin{gathered}
 F(t)=\frac{t(t+1)(2t+1)(3t+1)(4t+1)(5t+2)}{120},\quad
 \delta=t-1,\qquad w=\frac{F(t)}{4\delta^8},\\
 W=\prod_{i=1}^6\left(1+\frac{C_i^2}{F(t)}\right),\quad
 V=\Tr(W^{-1}),\quad b=\frac{4w}{V^2},\quad d=7b.
 \end{gathered}
\end{equation}
The detector product and its physical weight are unchanged.
Keep every original point, difference, quotient, product and trace
denominator guard.
For each block retain two rational parameters $v_j,s_j$ and adjoin
\begin{equation}\label{eq:pole}
 Y_j^2=P_{d_j}(v_j,s_j),\qquad
 P_d(v,s)=d(d+v^2)^2s^4-d(d-v^2)^2(1-d)^2.
\end{equation}
This equation needs no nonzero guards for pole soundness. One may
add the point-algebra norm of each cleared radicand to the encoder's
guard, ensuring that the final quadratic blocks are etale.

\section{Local extraction and soundness}
Let $S$ be the unchanged global exceptional set of the conditional
source; it contains two and seven and all fixed contact collisions.
Let $M_j$ be its original nonexceptional odd contact radical at $t_j$.

\begin{lemma}\label{lem:odd}
Every odd valuation of $w(t)$ outside $S$ is positive. At such a
valuation the intrinsic pole equation forces all
six tests to have positive valuation at a place above $p$.
\end{lemma}
\begin{proof}
If $v_p(t)<0$, then $v_p(F)=6v_p(t)$ and $v_p(\delta)=v_p(t)$,
so $v_p(w)=-2v_p(t)$ is even. If $v_p(\delta)>0$, then
$F(1)=7$ is a unit and $v_p(w)=-8v_p(\delta)$ is even.
Otherwise $w$ is integral. At an odd valuation $\nu=v_p(w)>0$,
the normalizing factor $4\delta^8$ is a unit, so $\nu=v_p(F)$.

Each detector factor has valuation $\min(0,2v(C_i)-\nu)\le0$,
with a unique minimum. Let $q\le0$ be the largest local valuation
of $W$ over $p$. Then $v_p(V)\ge-q$. By the unguarded local
theorem of \cite{Intrinsic}, $d$ has no negative odd valuation.
Thus $v_p(d)=\nu-2v_p(V)$ is positive and
$1\le v_p(d)\le\nu+2q$. At a place attaining $q$, each factor
valuation is at least $q>-\nu/2$, giving $v(C_i)>\nu/4>0$.
\end{proof}

\begin{theorem}\label{thm:main}
Under the retained group, prime-pattern and trace-height inputs, the
stated weak $abc$ inequality, and Burungale's Theorem B, the modified
source has integer input projection for every sufficiently large
fixed odd $B$. It is complete for every nonzero integer input and
every odd $B\ge3$. Including zero by the final input factor gives a
defining polynomial with thirteen total variables. No numerical
sufficient shift or complete determinant-degree bound is asserted.
\end{theorem}
\begin{proof}[Soundness]
The inverse transformation from $(C_5,C_6)$ to
$(r-\bar s,r-X)$ is invertible outside $S$. Lemma~\ref{lem:odd}
therefore supplies exactly the local conditions of the original
conditional argument. For $A=[n]P$, the formal comparison gives
$X\equiv n\pmod p$ at every consumed contact. If $X\ne n$, the
reduced numerator of $X-n$ is divisible by each $M_j$.
The arguments, contacts, scales and root point have not changed.
Hence the contact-height, binary-power and trace-height estimates and
the sufficient-shift contradiction of \cite{Conditional} apply with
the same constants and exponents. Multiplication of the pole input
by seven changes no valuation outside the retained exceptional set.
\end{proof}

\section{Ordered completeness}
Fix $X=n\ne0$ and odd $B\ge3$. Take $A=[n]P$ and fix both scales.
Temporarily set $H=[h]P$ and $U=[n]H$. Choose the original compact
real neighborhoods and fixed finite-place neighborhoods, with all
point and difference functions defined. The finitely many sign
branches may be included in the real avoidance and bounds.
The temporary finite preparation set includes seven, the original
primes two, three and seventeen, and all denominator and avoidance
primes. This does not enlarge the global exceptional set $S$.

At any fixed finite place, as $t$ becomes large,
\begin{equation}\label{eq:limits}
 F(t)\sim t^6,\quad w(t)\sim\frac1{4t^2},\quad
 W=1+O(t^{-6}),\quad V\longrightarrow2.
\end{equation}
The limits are uniform on the already chosen neighborhoods. At two,
take $t$ sufficiently large that $w$ is a local square with
$v_2(w)\ge4$. Then $v_2(V)=1$, so
\[
 b\in\Q_2^{\times2},\quad v_2(b)=v_2(w)\ge4,\quad
 d=7b\in7\Q_2^{\times2}.
\]
At seven make $w$ a local square of positive valuation and $V$ a unit.
Then $v_7(b)=v_7(w)\ge0$ and $v_7(d)>0$. The stronger original
preparation at three and seventeen may be kept.

At the real places take $t$ sufficiently large that $t>1$,
$0<w<1/28$ and $1\le W<2$ in both selected slots.
Then $V>1$ and $0<b<1/7$, whence $0<d<1$.
All finite depths and lower bounds on the slope are prescribed before
prime selection.

Apply the same simultaneous prime-pattern theorem to the unchanged
rational contact coefficients. Its fixed finite avoidance and
arbitrary denominator divisibility allow the added preparation at
seven. It supplies the common slope, prime contacts, and conjugate
reduction-order conditions. Refine $h$ in the same compatible
progression, preserving the chosen finite and real neighborhoods.

At a fresh selected contact, $v_p(w)=1$. At the successful
place, the first four tests have valuation at least two and the last
two at least four; each factor of $W$ is one modulo $p^3$. At the
other place its reciprocal product reduces to zero or one. Hence
$V$ is a unit and $v_p(b)=1$.
The original prime-times-square factorization shows that every other
possible odd valuation of $w$ belongs to the fixed preparation set.
There \eqref{eq:limits} supplies nonnegative valuations. All remaining
valuations of $w$, and hence of $b$, are even. Thus $b$ has
square reduced denominator. Since $v_7(b)\ge0$, so does $d=7b$.

Both actual inputs satisfy the intrinsic theorem's hypotheses:
\[
 0<d_j<1,\quad d_j\text{ has square reduced denominator},\quad
 d_j\in7\Q_2^{\times2},\quad v_2(d_j)\ge4.
\]
That theorem supplies independent rational pole witnesses.

\section{Every character and the finite encoding}
All original point and difference guards retain the source's signed-index
proof. The quantities $t_j,w_j$ depend on retained abscissas and
scalar coordinates, not on finite ordinate signs. At a prepared tuple,
$t_j>1$ and $w_j>0$ at every character.
Every point-algebra character is real, and every detector factor is
one plus a nonnegative square divided by $F(t_j)>0$. Thus all product,
reciprocal-trace and detector-input norms are units. In particular
$d_j$ is a unit, regardless of whether its other character values are
below one.

At a character with value $d_j=c\ne0$, $P_c(v,s)$ has nonzero leading
coefficient $c(c+v^2)^2$ in $s$. Its point-algebra norm is a nonzero
polynomial in $v,s$. The intrinsic theorem gives a Zariski-dense set
of actual witness pairs, so each pair can avoid this norm and every
additional nonzero polynomial guard. This makes the final quadratic
blocks etale without adding a free coordinate.

The three point blocks have rank $4^3=64$, and the two quadratic pole
blocks give total rank $256$. The one-coordinate guarded encoder of
\cite{Unit} applies. The unrestricted coordinates are the input, six
point coordinates, one slope, four pole parameters and one encoder
coordinate: thirteen in total. The final factor $X$ includes zero
without an additional coordinate. This proves Theorem~\ref{thm:main}.

\begin{thebibliography}{9}\raggedright
\bibitem{Conditional}\emph{Thirteen Variables from Direct Elliptic Index
Comparison}, repository manuscript,
\nolinkurl{papers/conditional_few_variables/direct_input13_20261110_v22.tex}.
\bibitem{Reweighted}\emph{Quadratic reweightings of native elliptic
detectors}, repository manuscript,
\nolinkurl{papers/pole_parameters/reweighted_elliptic_detectors_20261116_v53.tex}.
\bibitem{Unit}\emph{A single norm determinant for the 25- and 13-variable
constructions}, repository manuscript,
\nolinkurl{papers/finite_algebra/unit_norm_20261114_v28.tex}.
\bibitem{Intrinsic}\emph{A two-coordinate pole test of coefficient height
four}, accompanying manuscript, using A. Burungale,
\emph{A Diophantine definition of $\Z$ in $\Q$},
arXiv:2611.16130v1, Theorem B.
\end{thebibliography}
\end{document}
